\documentclass{beamer}
%\usepackage[bars]{beamerthemetree} % Beamer主题样式v 2.2
\usepackage{ctex}
\usepackage{bm}
\usepackage{cite}
\usepackage{color}
\usepackage{xcolor}
\usepackage{graphicx}
\usepackage{diagbox}

%\setbeamertemplate{bibliography item}[text]
\setbeamertemplate{bibliography item}{\insertbiblabel}
\setbeamertemplate{bibliography entry title}{}
\setbeamertemplate{bibliography entry location}{}
\setbeamertemplate{bibliography entry note}{}

\setbeamertemplate{caption}[numbered]


\usetheme{Warsaw} % Beamer主题样式v 3.0
% \usecolortheme{lily} % Beamer颜色主题样式
\usefonttheme{serif}

\graphicspath{{fig/}}

\newcommand{\upcite}[1]{\textsuperscript{\cite{#1}}}  %自定义命令\upcite, 使参考文献引用以上标出现
\newcommand{\M}[1]{\mathrm{#1}}

\newcommand{\R}[1]{\textcolor{red}{#1}}

\newcommand{\PP}[2]{\frac{\partial #1}{\partial #2}}
\newcommand{\PPP}[2]{\frac{\partial^{2} #1}{{\partial #2}^{2}}}

\newcommand{\DD}[2]{\frac{\mathrm{d} #1}{\mathrm{d} #2}}
\newcommand{\DDD}[2]{\frac{\mathrm{d}^{2} #1}{{\mathrm{d} #2}^{2}}}

    
\logo{\includegraphics[height=0.1\textwidth]{bnu-fig-logo.pdf}}
\title{等温捷径与绝热捷径构成的类卡诺热机的效率与功率}
\author{龚政楠}
\institute{北京师范大学}
\begin{document}

\begin{frame} % Cover slide
\titlepage
\end{frame}

\begin{frame}           %生成目录页，目录太长时加选项[shrink]
    %	\setcounter{page}{0}	%setcounter似乎对beamer无效
        \addtocounter{framenumber}{-2}%---------位置放在beginframe之后，不然无效
        \frametitle{目录}
        \tableofcontents 
\end{frame}

%也可以用 \frame{\titlepage}}代替 (在Beamer v 2.2 宏包中)
\section{引言：热机的功率与效率问题}

% \begin{frame}{引言：热机的功率与效率问题}
% \end{frame}

\begin{frame}{引言：热机的功率与效率问题}
    \begin{itemize}
        \item<2-> 1824，卡诺指出\upcite{2005}：工作在相同高温热源和相同低温热源的所有热机中，以可逆热机的效率最高，为$\eta _{\mathrm{C}}=1-\frac{\mathrm{c}}{\mathrm{h}}$ 
        \item<3-> 1975年，Curzon和Ahlborn研究了内可逆热机\upcite{Curzon1975}，得到了该热机在最大功率下的效率为$\eta _{\mathrm{CA}}=1-\sqrt{1-\eta _{\mathrm{C}}}$
        \item<4-> 2008年，T. Schmiedl 和 U. Seifert提出了布朗随机热机模型\upcite{Schmiedl2008}，最大功率下的效率为$\frac{2 \eta_{\mathrm{C}}}{4-\eta_{\mathrm{C}}}$
        \item<5-> 同年，涂展春推导出费曼棘轮热机\upcite{Tu2008}在最大功率下的效率$\frac{\eta _{\mathrm{C}}^{2}}{\eta _{\mathrm{C}}-\left( 1-\eta _{\mathrm{C}} \right) \ln \left( 1-\eta _{\mathrm{C}} \right)}$. \upcite{Tu2020}
        \item<6-> 2013年，涂展春在Schmiedl和Seifert工作\upcite{Schmiedl2008}的基础上，考虑了他们在过阻尼情况下忽略的惯性的影响，构建了一种类卡洛热机\upcite{Tu2013}. 发现这种随机热机在最大功率下的效率等于$\eta _{\mathrm{CA}}$
    \end{itemize}    
\end{frame}

\section{有限时间热力学}
\subsection{有限时间热机模型}
\begin{frame}{CA内可逆热机模型}
\begin{columns}
    \begin{column}{.62\textwidth}
        \begin{exampleblock}{Curzon和Ahlborn的假设}
            \begin{itemize}
                \item<2-> 内可逆假设
                \item<3-> 绝热过程熵产生为零
                \item<4-> 绝热过程与等温过程所用的时间成正比，即循环总时间为$t_0 = a (t_{\mathrm{h}}+t_{\mathrm{c}})$
                \item<5-> 线性热传输假设，即
                \begin{equation}
                    \begin{split}
                        Q_{\mathrm{h}}&=\kappa_{\mathrm{h}}\left(T_{\mathrm{h}}-T_{\mathrm{eh}}\right) t_{\mathrm{h}}\\
                        Q_{\mathrm{c}}&=\kappa_{\mathrm{c}}\left(T_{\mathrm{ec}}-T_{\mathrm{c}}\right) t_{\mathrm{c}}
                    \end{split} 
                    \label{eq1}
                \end{equation}
            \end{itemize}   
        \end{exampleblock}
    \end{column}

    \begin{column}{.38\textwidth}
        \begin{figure}[!htbp]
            \begin{center}
                \includegraphics[width=\textwidth]{p1.png}
            \end{center}
            \caption{CA内可逆热机}
            \label{f1}
        \end{figure}
    \end{column}
\end{columns}
\end{frame}

\begin{frame}{CA内可逆热机模型}
    \begin{columns}
        \begin{column}{.62\textwidth}
            \begin{block}{CA热机的功率}
                \pause
                输出功率
                \begin{equation}
                    P\left(T_{\mathrm{eh}}, T_{\mathrm{ec}}\right)=\frac{T_{\mathrm{eh}}-T_{\mathrm{ec}}}{\frac{T_{\mathrm{eh}}}{a \kappa_{h}\left(T_{h}-T_{\mathrm{eh}}\right)}+\frac{T_{\mathrm{ec}}}{a \kappa_{c}\left(T_{\mathrm{ec}}-T_{c}\right)}}
                    \label{eq2}
                \end{equation}
                \pause 针对$T_{\mathrm{eh}}, T_{\mathrm{ec}}$使功率取最大值
                \begin{equation}
                    \begin{split}
                        T_{\mathrm{eh}}^{*}&=\frac{\sqrt{\kappa_{h} T_{h}}+\sqrt{\kappa_{c} T_{c}}}{\sqrt{\kappa_{h}}+\sqrt{K_{c}}} \sqrt{T_{h}}\\
                        T_{\mathrm{ec}}^{*}&=\frac{\sqrt{\kappa_{h} T_{h}}+\sqrt{\kappa_{c} T_{c}}}{\sqrt{K_{h}}+\sqrt{\kappa_{c}}} \sqrt{T_{c}}
                    \end{split}
                    \label{eq3}
                \end{equation}
            \end{block}
        \end{column}
    
        \begin{column}{.38\textwidth}
            \onslide
            \begin{figure}
                \begin{center}
                    \includegraphics[width=\textwidth]{p1.png}
                \end{center}
                \setcounter{figure}{0}
                \caption{CA内可逆热机}
            \end{figure}
        \end{column}
    \end{columns}
\end{frame}

\begin{frame}{CA内可逆热机模型}
    \begin{alertblock}{CA热机在最大功率下的效率}
        根据$\eta=1-T_{\mathrm{ec}} / T_{\mathrm{eh}}$就可以得到
        \begin{equation}
            \eta _{\mathrm{CA}}=1-\sqrt{1-\eta _{\mathrm{C}}}
            \label{eq4}
        \end{equation}
    \end{alertblock}
\end{frame}

\begin{frame}{TU布朗随机热机模型}
    \begin{columns}
        \begin{column}{.62\textwidth}
            \begin{exampleblock}{TU布朗随机热机的假设}
                \begin{itemize}
                    \item<2->  两个绝热过程是在瞬间完成的，$t_0 = (t_{\mathrm{h}}+t_{\mathrm{c}})$
                    \item<3->  两个绝热过程中不可逆熵产生为零
                    \item<4->  过阻尼假设，不考虑惯性的影响
                \end{itemize}         
            \end{exampleblock}
        \end{column}
    
        \begin{column}{.38\textwidth}
            \onslide
            \begin{figure}
                \begin{center}
                    \includegraphics[width=\textwidth]{p2.png}
                \end{center}
                \caption{TU布朗随机热机模型}
                \label{f2}
            \end{figure}
        \end{column}
    \end{columns}
\end{frame}

\begin{frame}{TU布朗随机热机的功率}
    \begin{columns}
        \begin{column}{.62\textwidth}
            \begin{block}{TU布朗随机热机}
                \pause
                平均功输出  
                    \begin{equation}
                        \begin{split}
                            -W=\left(T_{\mathrm{h}}-T_{\mathrm{c}}\right) \Delta S-\left({1}/{t_{\mathrm{h}}}+{1}/{t_{\mathrm{c}}}\right) A_{\mathrm{irr}}
                        \end{split} 
                        \label{eq5}
                    \end{equation}
                    \pause
                    吸收热
                    \begin{equation}
                        \begin{split}
                            Q =T_{\mathrm{h}} \Delta S-{A_{\mathrm{irr}}}/{t_{\mathrm{h}}}
                        \end{split}
                        \label{eq6}
                    \end{equation}
                    \pause
                    功率
                    \begin{equation}
                        \begin{split}
                            P=\frac{\left(T_{\mathrm{h}}-T_{\mathrm{c}}\right) \Delta S-\left({1}/{t_{\mathrm{h}}}+{1}/{t_{\mathrm{c}}}\right) A_{\mathrm{irr}}}{t_{\mathrm{h}} + t_{\mathrm{c}}}
                        \end{split}
                        \label{eq7}
                    \end{equation}       
            \end{block}
        \end{column}
    
        \begin{column}{.38\textwidth}
            \onslide
            \begin{figure}
                \begin{center}
                    \includegraphics[width=\textwidth]{p2.png}
                \end{center}
                \setcounter{figure}{1}
                \caption{TU布朗随机热机模型}
            \end{figure}
        \end{column}
    \end{columns}
\end{frame}

\begin{frame}{TU布朗随机热机}
\begin{alertblock}{TU布朗随机热机的效率}
通过优化$t_1,\ t_2$不难得到，最大功率时的效率为
\begin{equation}
    \eta^{*}=\frac{\eta_{\mathrm{C}}}{2-\eta_{\mathrm{C}} / 2}
    \label{eq8}
\end{equation}
\end{alertblock}
\end{frame}

\subsection{绝热捷径}
\begin{frame}{绝热捷径}
为了在有限时间内实现平衡态的转化， Demirplak and Rice \upcite{Demirplak2003}和 Berry \upcite{Berry2009} 独立发展出了绝热捷径的策略。
\pause

依赖于时间缓变的$H_0(t)$，绝热定理\upcite{Griffiths2018}
\begin{equation}
\begin{split}
    |\psi(t)\rangle=&|n(t)\rangle e^{i \theta_{n}(t)} e^{i \gamma_{n}(t)},\\
    \theta_{n}(t)=-\frac{1}{\hbar} \int_{0}^{t} E_{n}(t^{\prime}) \M{d} t^{\prime},& \quad \gamma_{n}(t)=i \int_{0}^{t}\langle n(t^{\prime}) | \partial_{t^{\prime}}n(t^{\prime})\rangle \M{d} t^{\prime}            
\end{split}        
    \label{eq9}
\end{equation}
\pause
若有$H(t)$使得绝热定理严格成立
\pause
\begin{equation}
\begin{split}
    H(t)&=H_{0}(t)+i \hbar \sum_{m}\left(\left|\partial_{t} m\right\rangle\left\langle m\left|-\left\langle m \mid \partial_{t} m\right\rangle\right| m\right\rangle\langle m|\right) \\
    &\equiv H_0 + H_1
\end{split}
    \label{eq10}
\end{equation}
\end{frame}

\begin{frame}{绝热捷径}
    在经典的情形下，对于随时间缓变的$H_0 (t)$，也可以引入辅助哈密顿量$H_1 (t)$，使得绝热不变量严格守恒
        \begin{equation}
            I(E, \boldsymbol{\lambda}) \equiv \int \mathrm{d} \bm{\eta} \theta\left[E-H_{0}(z ; \boldsymbol{\lambda})\right]
          \label{eq11}
        \end{equation}
        \pause
    \begin{alertblock}{谐振子势场中的辅助势}
        根据Jarzynski的计算\upcite{Jarzynski2013}，对于依赖于时间的谐振子势
        \begin{equation}
            H_{0}(\bm{\eta} ; \lambda(t))=\frac{p^{2}}{2}+\frac{q^2}{2\lambda(t)^{2}}
            \label{eq12}
        \end{equation}
        \pause
        其相应的辅助势为
        \begin{equation}
            H_{1}(\bm{\eta} ; \lambda(t))=\frac{\dot{\lambda}}{2 \lambda} q p
            \label{eq13}
        \end{equation}
    \end{alertblock}
\end{frame}

\begin{frame}{利用绝热捷径实现平衡态的转化}
考虑谐振子势式\eqref{eq12}，施加辅助哈密顿量之后的新哈密顿量为
\begin{equation}
    H(t)=\frac{p^{2}}{2}+\frac{q^2}{2 \lambda(t)^{2}}+\frac{\dot{\lambda}}{2 \lambda(t)} q p,
    \label{eq13.5}
\end{equation}
\pause
可以得到运动方程
\begin{equation}
    \left\{\begin{aligned}
        &\dot{q}=\frac{\partial H}{\partial p}=\frac{p}{m}+\frac{\lambda(t)}{2 \lambda(t)} q \\
        &\dot{p}=-\frac{\partial H}{\partial q}=-2  \frac{q}{\lambda^{2}(t)}-\frac{\dot{\lambda}(t)}{2 \lambda(t)} p
    \end{aligned}\right.
    \label{eq14}
\end{equation}
\end{frame}

\begin{frame}{利用绝热捷径实现平衡态的转化}
\begin{alertblock}{平衡态的转化}
设初态为平衡态
\begin{equation}
    \rho_{i}=\frac{\beta_{i}}{2 \pi \lambda_{i}} \exp \left[-\beta_{i} H_0 \left(t_{i}\right)\right]
    \label{eq16}
\end{equation}
\pause
假设末态为等效温度为$\beta_{f}$的平衡态
\begin{equation}
    \rho_{f}=\frac{\beta_{f} }{2 \pi \lambda_{f}} \exp \left[-\beta_{f} H_0 \left(t_{f}\right)\right]
    \label{eq17}
\end{equation}
\pause
鉴于$\rho_{f}=\rho_{i}$，$\DD{H_0 (t)}{t}=-H_0 (t) \DD{\ln{\lambda(t)}}{t}$，得到\upcite{Tu2013}
\begin{equation}
    \frac{\beta_f}{\lambda_f} = \frac{\beta_i}{\lambda_i}
    \label{eq18}
\end{equation}
\end{alertblock}
\end{frame}
\subsection{等温捷径}
\begin{frame}{等温捷径}
为了实现等温平衡态之间的相互转化，类似于绝热捷径，李耿等人\upcite{Li2016}发展出了等温捷径的策略。

\pause

为系统施加一个额外的辅助势场，使得哈密顿量成为
\begin{equation}
    H(q, p, t)=H_{0}(q, p, \lambda(t))+U_{1}(q, p, t)
    \label{eq19}
\end{equation}
\pause
分布函数$\rho(q, p, t)$的演化由福克-普朗克方程决定
\begin{equation}
    \begin{split}
        \PP{\rho}{t} &+ \nabla \cdot \bm{J} = 0\\
        \bm{J} \equiv \rho \left( p + \PP{U}{p} \right) \hat{\mathbf{q}}&-\rho\left(\gamma p+\frac{\partial U}{\partial q}+ \PP{U}{p} + \frac{\gamma}{\rho \beta} \frac{\partial \rho}{\partial p}\right) \hat{\mathbf{p}}
    \end{split}
    \label{eq20}
\end{equation}
\end{frame}

\begin{frame}{等温捷径}
\begin{alertblock}{等温平衡态的转化}
系统在原哈密蹲下的瞬时平衡态
\begin{equation}
    \begin{split}
        \rho(q, p, t) &= \mathrm{e}^{\beta\left[F(\lambda(t))-H_{0}(x, p, \lambda(t))\right]}
    \end{split}
    \label{eq21}
\end{equation}
\pause
要使系统始终保持瞬时平衡态，将上式\eqref{eq21}代入方程\eqref{eq20}，得到\upcite{Li2016}
\begin{equation}
    \frac{\gamma}{\beta} \frac{\partial^{2} U_{1}}{\partial p^{2}}-\gamma p \frac{\partial U_{1}}{\partial p}+\frac{\partial U_{0}}{\partial q} \frac{\partial U_{1}}{\partial p}-p \frac{\partial U_{1}}{\partial q}=\left(\frac{d F}{d \lambda}-\frac{\partial U_{0}}{\partial \lambda}\right) \dot{\lambda}
    \label{eq22}
\end{equation}
\pause
利用上式\eqref{eq22}也可以得到，对谐振子势$U_0 (q, \lambda(t))= q^{2} / 2 \lambda(t)^2$
\begin{equation}
    U_{1}(q, p, t)=-\frac{\dot{\lambda}(t)}{2 \gamma \lambda(t)}\left[(p-\gamma q)^{2}+ q^{2}/\lambda(t)^2\right]
    \label{eq2.62}
\end{equation}
\end{alertblock}
\end{frame} 

\section{类卡诺热机的功率与效率}
\subsection{类卡诺热机模型}
\begin{frame}{类卡诺热机模型}
\begin{columns}
    \begin{column}{.6\textwidth}
        \begin{exampleblock}{类卡诺热机模型的假设}
        \begin{itemize}
        \item<2-> 忽略绝热捷径过程的时间 
        \end{itemize}
        \end{exampleblock}          
    \end{column}
    \begin{column}{.5\textwidth}
    \begin{figure}
        \centering
        \def\svgwidth{\columnwidth}
        \input{fig/p3.1.pdf_tex}
        \caption{类卡诺热机模型}
        \label{p3.1}
        \end{figure}
    \end{column}   
\end{columns}
\end{frame}

\begin{frame}{类卡诺热机模型}
    \begin{columns}
        \begin{column}{.6\textwidth}
            \begin{exampleblock}{类卡诺热机模型的假设}
            \begin{itemize}
            \item<2-> 忽略绝热捷径过程的时间 
            \end{itemize}
            \end{exampleblock}          
        \end{column}
        \begin{column}{.5\textwidth}
        \begin{figure}
            \centering
            \def\svgwidth{\columnwidth}
            \input{fig/p3.1.pdf_tex}
            \caption{类卡诺热机模型}
            \label{p3.1}
            \end{figure}
        \end{column}   
    \end{columns}
\end{frame}


\subsection{过阻尼布朗粒子}

\subsection{欠阻尼布朗粒子}


\footnotesize
\bibliography{ref/refs}
\bibliographystyle{unsrt}

\end{document}